\documentclass[a4paper,11pt]{article}
\usepackage[margin=27mm]{geometry}
\usepackage{fontspec,amsmath,amssymb,amsthm,unicode-math}
\setmainfont{STIXGeneral}\setmathfont{STIX Two Math}
\setlength{\emergencystretch}{2em}
\newcommand{\gen}[1]{\langle#1\rangle}
\newtheorem{theorem}{Theorem}\newtheorem{lemma}[theorem]{Lemma}\newtheorem{corollary}[theorem]{Corollary}
\title{Contraction of string markings through a central factor}
\author{Human Compression project}\date{Private draft; independent review pending}
\begin{document}\maketitle
\begin{abstract}
A commuting endpoint can be multiplied into the next surviving generator.
The resulting tail is a central twist of the old tail; its intersection
property follows directly from projection. When the lost endpoint can be
recovered in the first two tail generators, the standard end criterion
then gives a contraction of the entire marking. We give complete proofs,
including the rank-three obstruction and the conditions for iteration.
The contraction is the rank-reduction theorem of Brooksbank and Leemans;
we make its common central-factor argument explicit, without claiming
historical priority for the theorem or its consequences.
\end{abstract}

\section{Conventions and the two elementary ingredients}
A string C-marking $(s_0,\ldots,s_{n-1})$ consists of distinct nonidentity
involutions with $s_i s_j=s_j s_i$ when $|i-j|>1$, and
\[
 \gen{s_i:i\in I}\cap\gen{s_i:i\in J}=\gen{s_i:i\in I\cap J}
 \qquad(I,J\subseteq\{0,\ldots,n-1\}).
\]
It is irreducible when no adjacent pair commutes. Groups need not be finite.
We use the standard end criterion: a string marking is C if its two end
deletions are C and their groups intersect in the middle group. This is
McMullen--Schulte, Proposition 2E16, restated in [BL, Proposition 2.1].
All three premises will be checked at each use.

\begin{lemma}[Forward central twist]\label{central}
Let $(s_0,\ldots,s_{k-1})$ be a string C-marking of $H$, and let $z$ be
an involution in an independent direct factor. Replacing $s_0$ by $zs_0$
in $H\times\gen z$ gives a string C-marking of the subgroup generated
by the new tuple. The same holds at the other endpoint.
\end{lemma}
\begin{proof}
Projection onto $H$ sends the new marks to the old ones, so nonidentity,
distinctness and the string commutations hold. Induct on $k$, the rank-one
case being immediate. One end deletion is the old untwisted tail. The
other is a central twist of a shorter old marking, so is C by induction.
An element in the intersection projects into the old middle subgroup.
Since it lies in the untwisted end, it equals its projection. Every middle
generator is unchanged, so the intersection is exactly the new middle
subgroup. The end criterion finishes the induction. Reverse the tuple
for a twist at the other endpoint.
\end{proof}
This is the valid forward endpoint extension of Fernandes--Leemans--Mixer,
Proposition 5.3. Its converse is neither stated nor used. The same complete
argument occurs in our separate work on restriction kernels; no additional
result from that paper is a premise here.

\begin{lemma}[The first three marks]\label{three}
Suppose $(a,b,c)$ is an irreducible string C-marking. Then $b$ does not
commute with $ac$, and $\gen{a,b,c}$ is not dihedral.
\end{lemma}
\begin{proof}
The intersection property gives
$\gen{a,b}\cap\gen{b,c}=\gen b$ and makes $\gen{a,c}$ a Klein four
group. If $b$ commuted with $ac$, the central involution $ac$ would give
\[
 (ab)^2=(cb)^2\in\gen{a,b}\cap\gen{b,c}=\gen b.
\]
The square is not $1$, since $a,b$ do not commute. Nor can it be $b$:
the equation $abab=b$ implies $aba=1$ and hence $b=1$. This contradiction
proves the first assertion.

In a nonabelian dihedral group, every noncentral involution is a reflection.
None of $a,b,c$ can be central, because each has a noncommuting neighbour.
Thus all three would be reflections. Two distinct commuting reflections
in a finite dihedral group differ by its central half-turn; in an infinite
dihedral group they cannot commute. In the finite case $ac$ would therefore
be central, contradicting the first assertion. To see the assertion about
reflections directly, write them as $r^i t,r^j t$ with $trt=r^{-1}$.
Commutativity says $r^{2(i-j)}=1$, so their nontrivial product has order two
and is central. This also proves the infinite case.
\end{proof}
In particular a finite dihedral group of order at least six has an
irreducible C-marking precisely in rank two: a longer irreducible marking
would contain a triple prohibited by Lemma~\ref{three}, while the usual two
reflections give rank two. This proves [BL, Lemma 2.2] without importing its
cited Suzuki-group classification.

\section{The contraction theorem}
\begin{theorem}[Brooksbank--Leemans contraction]\label{contraction}
Let $(r_0,\ldots,r_{n-1})$ be an irreducible string C-marking of $G$,
where $n\geq4$. If
\begin{equation}\label{recover}
 r_0\in\gen{r_0r_2,r_3},
\end{equation}
then
\[
 (r_1,r_0r_2,r_3,\ldots,r_{n-1})
\]
is an irreducible string C-marking of the same group $G$.
\end{theorem}
\begin{proof}
The new elements are involutions. The product $r_0r_2$ is nonidentity
because the old marking has the intersection property. It cannot equal
any other retained mark, since that would put an old mark in the subgroup
generated by the others. The string commutations follow from the old ones.
Condition~\eqref{recover} recovers $r_0$ and then $r_2$, so the generated
group is $G$.

For $n=4$ put $a=r_0,b=r_1,c=r_2,d=r_3$. Both new end pairs are C:
two distinct involutions generate a rank-two C-group. Their intersection
is contained in
\[
 \gen{a,b,c}\cap\gen{a,c,d}=\gen{a,c}.
\]
It contains $ac$. If it were larger than $\gen{ac}$, it would be the
whole Klein four group $\gen{a,c}$. In particular $a,c$ would lie in
$\gen{b,ac}$, making $\gen{a,b,c}=\gen{b,ac}$ dihedral, contrary to
Lemma~\ref{three}. Hence the new end intersection is exactly $\gen{ac}$.
The end criterion proves the base case. Notice that this part proves C for
the generated subgroup even without~\eqref{recover}; that condition is
needed to identify the subgroup with $G$.

Now let $n>4$ and assume the result for smaller ranks. The first end of the
new marking is the contraction of the old prefix through $r_{n-2}$.
That prefix is irreducible and contains all the elements in
\eqref{recover}, so induction makes this end C and identifies its group as
\[
 H=\gen{r_0,\ldots,r_{n-2}}.
\]
The other end is $(r_0r_2,r_3,\ldots,r_{n-1})$. The old intersection
property and string commutations give the actual direct product
\[
 \gen{r_0,r_2,\ldots,r_{n-1}}
 =\gen{r_0}\times\gen{r_2,\ldots,r_{n-1}}.
\]
Thus Lemma~\ref{central} makes the new end C. Condition~\eqref{recover}
identifies its group with the whole displayed product, denoted by $K$.
The old intersection property now gives
\[
 H\cap K=\gen{r_0,r_2,\ldots,r_{n-2}}
        =\gen{r_0r_2,r_3,\ldots,r_{n-2}},
\]
where the last equality again uses~\eqref{recover}. This is precisely the
new middle group, so the end criterion completes the induction.

Finally the new first adjacent pair does not commute by
Lemma~\ref{three}. If $r_0r_2$ commuted with $r_3$, then $r_2$ would
commute with $r_3$, since $r_0$ does. All other adjacent pairs are unchanged.
The new marking is therefore irreducible.
\end{proof}

\section{Parity and iteration}
\begin{lemma}\label{parity}
In a string C-marking, the condition~\eqref{recover} holds if and only if
$r_2r_3$ has finite odd order.
\end{lemma}
\begin{proof}
The old marked subgroup is
$\gen{r_0,r_2,r_3}=\gen{r_0}\times\gen{r_2,r_3}$.
Write $m=|r_2r_3|$. If $m$ is odd, then
$((r_0r_2)r_3)^m=r_0$, proving recovery. If $m$ is even or infinite,
the dihedral group $D=\gen{r_2,r_3}$ has a homomorphism
$\chi:D\to C_2$ taking $r_2$ to its nonidentity element and $r_3$ to $1$.
For finite $m$ the defining relation $(r_2r_3)^m=1$ respects these values
exactly because $m$ is even; for infinite $m$ there is no such relation.
The subgroup $\gen{r_0r_2,r_3}$ lies in the graph of $\chi$ inside
$\gen{r_0}\times D$, and that graph excludes $(r_0,1)$. This proves the
converse. This parity argument already appears in our earlier analysis of
rank-three folds; it is an explicit use of the restriction kernel, not a
necessary condition for every successful contraction.
\end{proof}

\begin{corollary}\label{iterate}
Let the original irreducible rank-$n$ marking have adjacent orders
$p_j=|r_{j-1}r_j|$. Suppose $0\leq t\leq n-3$ and
$p_3,p_4,\ldots,p_{t+2}$ are all finite and odd, with the condition empty
when $t=0$. Then the same group has irreducible string C-markings in every
rank $n,n-1,\ldots,n-t$.
\end{corollary}
\begin{proof}
For one contraction, Lemma~\ref{parity} and Theorem~\ref{contraction}
apply. The third adjacent product in the new tuple is $r_3r_4$, of order
$p_4$; subsequent adjacent products are also unchanged. Thus at the next
step the same test uses $p_4$, then $p_5$, and so on. Irreducibility is
preserved by the theorem at every step. Stop at rank three. This proves
each indicated rank by induction.
\end{proof}
The first two adjacent orders can change, but neither is the next test.
In particular it is unnecessary to recompute a full subgroup census after
each contraction. The hypotheses and proof are those of the published
rank-reduction method, here organized around the central-factor argument.
No claim of historical novelty or new group existence beyond its stated
applications is made.

\begin{thebibliography}{9}
\bibitem{BL} P. A. Brooksbank and D. Leemans, \emph{Rank reduction of string
C-group representations}, arXiv:1812.01055v2, 3 May 2019; published in
Proc. Amer. Math. Soc. 147 (2019), 5421--5426.
\bibitem{FLM} M. E. Fernandes, D. Leemans and M. Mixer, \emph{All alternating
groups $A_n$ with $n\geq12$ have polytopes of rank $\lfloor(n-1)/2\rfloor$},
SIAM J. Discrete Math. 26 (2012), 482--498, Proposition 5.3.
\end{thebibliography}
\end{document}
